\documentclass[11pt,a4paper,DIV=12]{scrartcl}
\usepackage{multicol}

\makeatletter
\def\input@path{{../../../}}
\makeatother

\usepackage[sexy]{evan}
% \usepackage[sexy]{../../../evan}

% Override environment "problem" agar hanya menampilkan angka
\makeatletter
\newtheoremstyle{justnumber}
  {8pt} % space above
  {8pt} % space below
  {\normalfont} % body font
  {} % indent amount
  {\color{blue!40!black}\bfseries} % head font
  {} % punctuation after head
  {0.5em} % space after head
  {\thmnumber{#2}.\thmnote{ (\small #3)}} % head spec (hanya angka dan note)

% Hapus definisi problem lama di dalam makeatletter agar @ terbaca
\let\problem\relax
\let\endproblem\relax
\let\c@problem\relax
\let\theproblem\relax
\makeatother

\theoremstyle{justnumber}
\newtheorem{problem}{}

\begin{document}
\title{Latihan Soal 2:
  Identitas, Barisan, Polinomial}
\author{MAN 1 Pasuruan}
\date{Mei 2025}
\maketitle
\section*{Pilihan Ganda}
\begin{problem}[AIME 1983]
  Tentukan hasil kali akar real dari persamaan
  \[
    x^2 + 18x + 30 = 2\sqrt{x^2 + 18x + 45}.
  \]
  \begin{multicols}{5}
    \begin{enumerate}
      \item[(A)] -20
      \item[(B)] 15
      \item[(C)] 20
      \item[(D)] 30
      \item[(E)] 45
    \end{enumerate}
  \end{multicols}
  %KJ: C
\end{problem}

\begin{problem}[AIME 1983]
  Diberikan barisan aritmetika $a_1, a_2, a_3, \ldots$ dengan beda $1$. Jika
  \[
    a_1 + a_2 + a_3 + \cdots + a_{98} = 137,
  \]
  tentukan nilai
  \[
    a_2 + a_4 + a_6 + \cdots + a_{98}.
  \]

  \begin{multicols}{5}
    \begin{enumerate}
      \item[(A)] 44
      \item[(B)] 68
      \item[(C)] 88
      \item[(D)] 93
      \item[(E)] 186
    \end{enumerate}
  \end{multicols}
  %KJ: D
\end{problem}

\begin{problem}[OSK 2017]
  Misalkan $P(x)$ suatu polinom berderajat $4$ yang memiliki nilai maksimum $2018$ di $x=0$ dan $x=2$. Jika $P(1)=2017$, maka nilai $P(3)$ adalah \ldots
  \begin{multicols}{5}
    \begin{enumerate}
      \item[(A)] 2000
      \item[(B)] 2009
      \item[(C)] 2015
      \item[(D)] 2018
      \item[(E)] 2027
    \end{enumerate}
  \end{multicols}
\end{problem}

\section*{Isian Singkat}
\begin{problem}[AIME II 2005]
  Misalkan
  \[
    x = \frac{4}{(\sqrt{5}+1)(\sqrt[4]{5}+1)(\sqrt[8]{5}+1)(\sqrt[16]{5}+1)}.
  \]
  Tentukan nilai dari $(x+1)^{48}$.
  %KJ: 125
\end{problem}

\begin{problem}[AIME II 2005]
  Misalkan $P(x)$ adalah polinomial dengan koefisien bilangan bulat yang memenuhi $P(17) = 10$ dan $P(24) = 17$. Jika persamaan $P(n) = n + 3$ memiliki dua solusi bilangan bulat yang berbeda $n_1$ dan $n_2$, tentukan hasil kali $n_1 \cdot n_2$.
  %KJ: 418
\end{problem}

\begin{problem}[OSK 2022]
  Misalkan $f(x) = a^2 x + 200$. Jika $f(20) + f^{-1}(22) = f^{-1}(20) + f(22)$, maka $f(1) = \ldots$
\end{problem}

\end{document}